% Options for packages loaded elsewhere
% Options for packages loaded elsewhere
\PassOptionsToPackage{unicode}{hyperref}
\PassOptionsToPackage{hyphens}{url}
\PassOptionsToPackage{dvipsnames,svgnames,x11names}{xcolor}
%
\documentclass[
  12pt,
  letterpaper]{article}
\usepackage{xcolor}
\usepackage[margin=1in]{geometry}
\usepackage{amsmath,amssymb}
\setcounter{secnumdepth}{5}
\usepackage{iftex}
\ifPDFTeX
  \usepackage[T1]{fontenc}
  \usepackage[utf8]{inputenc}
  \usepackage{textcomp} % provide euro and other symbols
\else % if luatex or xetex
  \usepackage{unicode-math} % this also loads fontspec
  \defaultfontfeatures{Scale=MatchLowercase}
  \defaultfontfeatures[\rmfamily]{Ligatures=TeX,Scale=1}
\fi
\usepackage{lmodern}
\ifPDFTeX\else
  % xetex/luatex font selection
\fi
% Use upquote if available, for straight quotes in verbatim environments
\IfFileExists{upquote.sty}{\usepackage{upquote}}{}
\IfFileExists{microtype.sty}{% use microtype if available
  \usepackage[]{microtype}
  \UseMicrotypeSet[protrusion]{basicmath} % disable protrusion for tt fonts
}{}
\usepackage{setspace}
\makeatletter
\@ifundefined{KOMAClassName}{% if non-KOMA class
  \IfFileExists{parskip.sty}{%
    \usepackage{parskip}
  }{% else
    \setlength{\parindent}{0pt}
    \setlength{\parskip}{6pt plus 2pt minus 1pt}}
}{% if KOMA class
  \KOMAoptions{parskip=half}}
\makeatother
% Make \paragraph and \subparagraph free-standing
\makeatletter
\ifx\paragraph\undefined\else
  \let\oldparagraph\paragraph
  \renewcommand{\paragraph}{
    \@ifstar
      \xxxParagraphStar
      \xxxParagraphNoStar
  }
  \newcommand{\xxxParagraphStar}[1]{\oldparagraph*{#1}\mbox{}}
  \newcommand{\xxxParagraphNoStar}[1]{\oldparagraph{#1}\mbox{}}
\fi
\ifx\subparagraph\undefined\else
  \let\oldsubparagraph\subparagraph
  \renewcommand{\subparagraph}{
    \@ifstar
      \xxxSubParagraphStar
      \xxxSubParagraphNoStar
  }
  \newcommand{\xxxSubParagraphStar}[1]{\oldsubparagraph*{#1}\mbox{}}
  \newcommand{\xxxSubParagraphNoStar}[1]{\oldsubparagraph{#1}\mbox{}}
\fi
\makeatother


\usepackage{longtable,booktabs,array}
\usepackage{calc} % for calculating minipage widths
% Correct order of tables after \paragraph or \subparagraph
\usepackage{etoolbox}
\makeatletter
\patchcmd\longtable{\par}{\if@noskipsec\mbox{}\fi\par}{}{}
\makeatother
% Allow footnotes in longtable head/foot
\IfFileExists{footnotehyper.sty}{\usepackage{footnotehyper}}{\usepackage{footnote}}
\makesavenoteenv{longtable}
\usepackage{graphicx}
\makeatletter
\newsavebox\pandoc@box
\newcommand*\pandocbounded[1]{% scales image to fit in text height/width
  \sbox\pandoc@box{#1}%
  \Gscale@div\@tempa{\textheight}{\dimexpr\ht\pandoc@box+\dp\pandoc@box\relax}%
  \Gscale@div\@tempb{\linewidth}{\wd\pandoc@box}%
  \ifdim\@tempb\p@<\@tempa\p@\let\@tempa\@tempb\fi% select the smaller of both
  \ifdim\@tempa\p@<\p@\scalebox{\@tempa}{\usebox\pandoc@box}%
  \else\usebox{\pandoc@box}%
  \fi%
}
% Set default figure placement to htbp
\def\fps@figure{htbp}
\makeatother


% definitions for citeproc citations
\NewDocumentCommand\citeproctext{}{}
\NewDocumentCommand\citeproc{mm}{%
  \begingroup\def\citeproctext{#2}\cite{#1}\endgroup}
\makeatletter
 % allow citations to break across lines
 \let\@cite@ofmt\@firstofone
 % avoid brackets around text for \cite:
 \def\@biblabel#1{}
 \def\@cite#1#2{{#1\if@tempswa , #2\fi}}
\makeatother
\newlength{\cslhangindent}
\setlength{\cslhangindent}{1.5em}
\newlength{\csllabelwidth}
\setlength{\csllabelwidth}{3em}
\newenvironment{CSLReferences}[2] % #1 hanging-indent, #2 entry-spacing
 {\begin{list}{}{%
  \setlength{\itemindent}{0pt}
  \setlength{\leftmargin}{0pt}
  \setlength{\parsep}{0pt}
  % turn on hanging indent if param 1 is 1
  \ifodd #1
   \setlength{\leftmargin}{\cslhangindent}
   \setlength{\itemindent}{-1\cslhangindent}
  \fi
  % set entry spacing
  \setlength{\itemsep}{#2\baselineskip}}}
 {\end{list}}
\usepackage{calc}
\newcommand{\CSLBlock}[1]{\hfill\break\parbox[t]{\linewidth}{\strut\ignorespaces#1\strut}}
\newcommand{\CSLLeftMargin}[1]{\parbox[t]{\csllabelwidth}{\strut#1\strut}}
\newcommand{\CSLRightInline}[1]{\parbox[t]{\linewidth - \csllabelwidth}{\strut#1\strut}}
\newcommand{\CSLIndent}[1]{\hspace{\cslhangindent}#1}



\setlength{\emergencystretch}{3em} % prevent overfull lines

\providecommand{\tightlist}{%
  \setlength{\itemsep}{0pt}\setlength{\parskip}{0pt}}



 


\usepackage{booktabs}
\usepackage{longtable}
\usepackage{array}
\usepackage{multirow}
\usepackage{wrapfig}
\usepackage{float}
\usepackage{colortbl}
\usepackage{pdflscape}
\usepackage{tabu}
\usepackage{threeparttable}
\usepackage{threeparttablex}
\usepackage[normalem]{ulem}
\usepackage{makecell}
\usepackage{xcolor}
\usepackage{setspace}
\usepackage{booktabs}
\usepackage{caption}
\usepackage{longtable}
\usepackage{array}
\usepackage{multirow}
\usepackage{float}
\usepackage{graphicx}
\usepackage{threeparttable}
\usepackage{threeparttablex}
\captionsetup{width=\textwidth}
\setlength\LTcapwidth{\textwidth}
\floatplacement{figure}{H}
\floatplacement{table}{H}
\raggedright
\makeatletter
\@ifpackageloaded{caption}{}{\usepackage{caption}}
\AtBeginDocument{%
\ifdefined\contentsname
  \renewcommand*\contentsname{Table of contents}
\else
  \newcommand\contentsname{Table of contents}
\fi
\ifdefined\listfigurename
  \renewcommand*\listfigurename{List of Figures}
\else
  \newcommand\listfigurename{List of Figures}
\fi
\ifdefined\listtablename
  \renewcommand*\listtablename{List of Tables}
\else
  \newcommand\listtablename{List of Tables}
\fi
\ifdefined\figurename
  \renewcommand*\figurename{Figure}
\else
  \newcommand\figurename{Figure}
\fi
\ifdefined\tablename
  \renewcommand*\tablename{Table}
\else
  \newcommand\tablename{Table}
\fi
}
\@ifpackageloaded{float}{}{\usepackage{float}}
\floatstyle{ruled}
\@ifundefined{c@chapter}{\newfloat{codelisting}{h}{lop}}{\newfloat{codelisting}{h}{lop}[chapter]}
\floatname{codelisting}{Listing}
\newcommand*\listoflistings{\listof{codelisting}{List of Listings}}
\makeatother
\makeatletter
\makeatother
\makeatletter
\@ifpackageloaded{caption}{}{\usepackage{caption}}
\@ifpackageloaded{subcaption}{}{\usepackage{subcaption}}
\makeatother
\usepackage{bookmark}
\IfFileExists{xurl.sty}{\usepackage{xurl}}{} % add URL line breaks if available
\urlstyle{same}
\hypersetup{
  pdftitle={Online Appendix: The Satisfaction Paradox},
  pdfauthor={Jeffrey Stark},
  colorlinks=true,
  linkcolor={blue},
  filecolor={Maroon},
  citecolor={blue},
  urlcolor={blue},
  pdfcreator={LaTeX via pandoc}}


\title{Online Appendix: The Satisfaction Paradox}
\usepackage{etoolbox}
\makeatletter
\providecommand{\subtitle}[1]{% add subtitle to \maketitle
  \apptocmd{\@title}{\par {\large #1 \par}}{}{}
}
\makeatother
\subtitle{Economic Performance and the Structure of Democratic Support
in South Korea}
\author{Jeffrey Stark}
\date{March 2, 2026}
\begin{document}
\maketitle


\setstretch{2}
\section{Appendix A: Data and
Measurement}\label{appendix-a-data-and-measurement}

\subsection{A.1 Survey Coverage}\label{a.1-survey-coverage}

The analysis draws on six waves of the Asian Barometer Survey (ABS), a
comparative survey program administering nationally representative
face-to-face interviews across East and Southeast Asia (Asian Barometer
Survey, 2023). Table~\ref{tbl-coverage} reports wave-level sample sizes
for Korea and Taiwan.

\begin{table}

\caption{\label{tbl-coverage}Sample coverage: Korea and Taiwan, ABS
waves 1--6}

\centering{

\centering
\begin{tabular}[t]{lrrrr}
\toprule
\multicolumn{1}{c}{ } & \multicolumn{2}{c}{Korea} & \multicolumn{2}{c}{Taiwan} \\
\cmidrule(l{3pt}r{3pt}){2-3} \cmidrule(l{3pt}r{3pt}){4-5}
Wave & Year & N & Year & N\\
\midrule
1 & 2,003 & 1,500 & 2,001 & 1,415\\
2 & 2,006 & 1,212 & 2,006 & 1,587\\
3 & 2,011 & 1,207 & 2,010 & 1,592\\
4 & 2,015 & 1,200 & 2,014 & 1,657\\
5 & 2,019 & 1,268 & 2,019 & 1,259\\
\addlinespace
6 & 2,022 & 1,216 & 2,022 & 1,532\\
\bottomrule
\end{tabular}

}

\end{table}%

\subsection{A.2 Variable Construction and Question
Wording}\label{a.2-variable-construction-and-question-wording}

All survey items were normalized to a 0--1 scale within country using
min--max normalization. Indices are row-wise means of their normalized
components.

\textbf{Satisfaction index} (2 items, all waves):

\begin{itemize}
\tightlist
\item
  \emph{Democracy satisfaction} (\texttt{democracy\_satisfaction}): ``On
  the whole, how satisfied or dissatisfied are you with the way
  democracy works in our country?'' (1 = very dissatisfied, 4 = very
  satisfied)
\item
  \emph{Government performance} (\texttt{gov\_sat\_national}): ``How
  satisfied are you with the performance of the current national
  government?'' (1 = very dissatisfied, 4 = very satisfied)
\end{itemize}

\textbf{Democratic quality index} (2 items, all waves):

\begin{itemize}
\tightlist
\item
  \emph{Always preferable} (\texttt{dem\_always\_preferable}): ``Which
  of the following statements comes closest to your own opinion? (a)
  Democracy is always preferable to any other kind of government; (b)
  Under some circumstances, an authoritarian government can be
  preferable; (c) For people like me, it doesn't matter whether we have
  a democratic or a non-democratic government.'' (Coded 1 = always
  preferable, 0 otherwise, then normalized.)
\item
  \emph{Extent democratic} (\texttt{dem\_extent\_current}): ``How much
  of a democracy is {[}country{]} today?'' (1--10 scale)
\end{itemize}

\textbf{Authoritarian rejection index} (4 items; higher = more rejection
of authoritarianism):

``I'm going to describe some political systems. For each, tell me
whether you think it would be very good, fairly good, fairly bad, or
very bad for governing {[}country{]}.''

\begin{itemize}
\tightlist
\item
  ``Having a strong leader who does not have to bother with parliament
  and elections'' (strongman rule)
\item
  ``Having the army rule the country'' (military rule)
\item
  ``Having experts, not government, make decisions according to what
  they think is best for the country'' (expert rule)
\item
  ``Having only one political party allowed to stand for election and
  rule the country'' (single-party rule)
\end{itemize}

Each item scored 1 = very good \ldots{} 4 = very bad. Items are
reverse-coded before normalization so that higher values indicate
stronger rejection of authoritarianism.

\textbf{Economic evaluation index} (6 items, all waves):

Six items spanning sociotropic and pocketbook assessments across
current, retrospective, and prospective time horizons: national economy
now; household economy now; national economic outlook; household
economic outlook; national economic change; household economic change.
All scored 1--5 (worse to better) and normalized 0--1.

\textbf{Controls} (Dalton, 2004): age (normalized 0--1), gender (1 =
male), education level (normalized 0--1), urban/rural residence (dummy),
political interest (normalized 0--1).

\subsection{A.3 Wave-Level Descriptive
Statistics}\label{a.3-wave-level-descriptive-statistics}

\begin{table}

\caption{\label{tbl-desc-kr}Korea: wave-level index means (normalized
0--1, SE in parentheses)}

\centering{

\centering
\resizebox{\ifdim\width>\linewidth\linewidth\else\width\fi}{!}{
\begin{tabular}[t]{lrrlll}
\toprule
Wave & Year & N & Satisfaction & Dem. quality & Econ. eval.\\
\midrule
1 & 2,003 & 1,500 & 0.479 (0.004) & 0.370 (0.008) & 0.443 (0.003)\\
2 & 2,006 & 1,212 & 0.386 (0.005) & 0.227 (0.005) & 0.356 (0.004)\\
3 & 2,011 & 1,207 & 0.457 (0.006) & 0.171 (0.005) & 0.423 (0.004)\\
4 & 2,015 & 1,200 & 0.511 (0.006) & 0.182 (0.005) & 0.423 (0.004)\\
5 & 2,019 & 1,268 & 0.509 (0.005) & 0.151 (0.004) & 0.382 (0.003)\\
\addlinespace
6 & 2,022 & 1,216 & 0.524 (0.005) & 0.141 (0.004) & 0.384 (0.003)\\
\bottomrule
\multicolumn{6}{l}{\textsuperscript{} SE = standard error of the mean.}\\
\end{tabular}}

}

\end{table}%

\begin{table}

\caption{\label{tbl-desc-tw}Taiwan: wave-level index means (normalized
0--1, SE in parentheses)}

\centering{

\centering
\begin{tabular}[t]{lrrrr}
\toprule
Wave & Year & N & Satisfaction & Dem. quality\\
\midrule
1 & 2,001 & 1,415 & 0.458 (0.005) & 0.431 (0.010)\\
2 & 2,006 & 1,587 & 0.431 (0.005) & 0.276 (0.006)\\
3 & 2,010 & 1,592 & 0.500 (0.005) & 0.245 (0.005)\\
4 & 2,014 & 1,657 & 0.436 (0.005) & 0.273 (0.006)\\
5 & 2,019 & 1,259 & 0.466 (0.006) & 0.265 (0.006)\\
\addlinespace
6 & 2,022 & 1,532 & 0.566 (0.006) & 0.238 (0.006)\\
\bottomrule
\multicolumn{5}{l}{\textsuperscript{} Econ. eval. index not shown for Taiwan as it is not the}\\
\multicolumn{5}{l}{focus of the paper.}\\
\end{tabular}

}

\end{table}%

\section{Appendix B: Full Wave-by-Wave OLS
Results}\label{appendix-b-full-wave-by-wave-ols-results}

Table~\ref{tbl-wave-all} presents OLS coefficients on
\texttt{econ\_index} for all five dependent variables across all six
waves, for Korea and Taiwan separately. All models include controls for
age, gender, education, urban/rural residence, and political interest.
The table underpins Figures 3 and 4 in the main text.

\begin{table}

\caption{\label{tbl-wave-all}Wave-by-wave OLS: econ to DV for all
outcomes, Korea and Taiwan}

\centering{

\centering\begingroup\fontsize{9}{11}\selectfont

\resizebox{\ifdim\width>\linewidth\linewidth\else\width\fi}{!}{
\begin{tabular}[t]{llrrrrr}
\toprule
Country & Year & Satisfaction with democracy & Satisfaction with government & Dem always preferable & Democratic extent & Auth rejection index\\
\midrule
 & 2003 & 0.278*** & 0.473*** & -0.246** & -0.124*** & -0.046\\

 & 2006 & 0.262*** & 0.522*** & -0.121 & 0.051*** & -0.034\\

 & 2011 & 0.361*** & 0.556*** & 0.067 & 0.068*** & 0.012\\

 & 2015 & 0.455*** & 0.474*** & 0.053 & 0.106*** & -0.224***\\

 & 2019 & 0.350*** & 0.823*** & 0.122 & 0.108*** & -0.161**\\

\multirow{-6}{*}{\raggedright\arraybackslash Korea} & 2022 & 0.339*** & 0.569*** & 0.079 & 0.071*** & -0.089\\
\cmidrule{1-7}
 & 2001 & 0.273*** & 0.491*** & -0.287*** & -0.110 & -0.015\\

 & 2006 & 0.443*** & 0.605*** & -0.306*** & 0.109*** & -0.065*\\

 & 2010 & 0.292*** & 0.634*** & 0.065 & 0.100*** & -0.112***\\

 & 2014 & 0.390*** & 0.577*** & -0.139* & 0.106*** & -0.088***\\

 & 2019 & 0.397*** & 0.733*** & -0.271*** & 0.133*** & -0.011\\

\multirow{-6}{*}{\raggedright\arraybackslash Taiwan} & 2022 & 0.561*** & 0.841*** & -0.504*** & 0.183*** & 0.074**\\
\bottomrule
\multicolumn{7}{l}{\textsuperscript{} OLS coefficient on econ\_index. Controls: age, gender, education, urban/rural, political interest. *p < .05, **p < .01, ***p < .001.}\\
\end{tabular}}
\endgroup{}

}

\end{table}%

\section{Appendix C: Pooled Individual-DV
Models}\label{appendix-c-pooled-individual-dv-models}

Table~\ref{tbl-pooled} reports pooled OLS results (with wave fixed
effects) for each individual dependent variable in Korea and Taiwan.
This complements the wave-by-wave table by presenting a single summary
estimate for each outcome across the full survey period.

\begin{table}

\caption{\label{tbl-pooled}Pooled OLS: econ\(\to\)individual DVs, Korea
and Taiwan}

\centering{

\centering
\begin{threeparttable}
\begin{tabular}[t]{>{\raggedright\arraybackslash}p{4.5cm}>{\raggedleft\arraybackslash}p{2.8cm}r>{\raggedleft\arraybackslash}p{2.8cm}r}
\toprule
\multicolumn{1}{c}{ } & \multicolumn{2}{c}{Korea} & \multicolumn{2}{c}{Taiwan} \\
\cmidrule(l{3pt}r{3pt}){2-3} \cmidrule(l{3pt}r{3pt}){4-5}
Dependent variable & $\beta$ (SE) & N & $\beta$ (SE) & N\\
\midrule
Satisfaction with democracy & 0.342*** (0.017) & 7,392 & 0.398*** (0.015) & 8,648\\
Satisfaction with government & 0.575*** (0.021) & 7,338 & 0.660*** (0.017) & 8,527\\
Dem always preferable & -0.032 (0.030) & 7,267 & -0.239*** (0.028) & 8,527\\
Democratic extent & 0.059*** (0.006) & 6,655 & 0.114*** (0.008) & 7,795\\
Auth rejection index & -0.082*** (0.018) & 7,519 & -0.040*** (0.011) & 8,729\\
\bottomrule
\end{tabular}
\begin{tablenotes}
\item OLS b (econ\_index) with wave FE and standard controls. SE in parentheses. *p < .05, **p < .01, ***p < .001.
\end{tablenotes}
\end{threeparttable}

}

\end{table}%

\section{Appendix D: Cross-Country Interaction
Test}\label{appendix-d-cross-country-interaction-test}

Table~\ref{tbl-cross-country} tests the H4 prediction that the
Korea--Taiwan divergence in the economic--quality relationship is
statistically significant. Taiwan is the reference category; a positive
\texttt{econ\_index:is\_korea} interaction indicates that the positive
economic effect on satisfaction is stronger in Korea than in Taiwan; a
negative coefficient indicates the reverse (i.e., Taiwan shows a
stronger negative effect on quality, as expected under the
critical-citizens pattern).

\begin{table}

\caption{\label{tbl-cross-country}Cross-country interaction: econ x
Korea, all five outcomes}

\centering{

\centering\begingroup\fontsize{9}{11}\selectfont

\resizebox{\ifdim\width>\linewidth\linewidth\else\width\fi}{!}{
\begin{tabular}[t]{lrrrrr}
\toprule
Term & Satisfaction with democracy & Satisfaction with government & Dem always preferable & Democratic extent & Auth rejection index\\
\midrule
Econ. eval. & 0.408*** (0.014) & 0.643*** (0.016) & -0.228*** (0.026) & 0.127*** (0.006) & -0.011 (0.012)\\
Korea (vs. Taiwan) & 0.043** (0.014) & -0.005 (0.016) & -0.008 (0.025) & 0.015* (0.006) & 0.007 (0.012)\\
Econ. x Korea & -0.090*** (0.023) & -0.064* (0.027) & 0.217*** (0.042) & -0.066*** (0.011) & -0.088*** (0.020)\\
\bottomrule
\multicolumn{6}{l}{\textsuperscript{} OLS. Taiwan = reference category. Pooled model with wave FE and standard controls. SE in parentheses. *p < .05, **p < .01, ***p < .001.}\\
\end{tabular}}
\endgroup{}

}

\end{table}%

\section{Appendix E: Subgroup
Heterogeneity}\label{appendix-e-subgroup-heterogeneity}

\subsection{E.1 Political Interest
Subgroups}\label{e.1-political-interest-subgroups}

H5 predicts the satisfaction--quality decoupling is not concentrated in
a single demographic group. Table~\ref{tbl-polint} presents results by
political interest tertile (low, medium, high) for Korea. All subgroups
should show positive econ → satisfaction and near-zero or negative econ
→ quality.

\begin{table}

\caption{\label{tbl-polint}Korea: econ to DV by political interest
subgroup}

\centering{

\centering
\begin{tabular}[t]{lrrr}
\toprule
Group & Dem always preferable & Satisfaction with democracy & Auth rejection index\\
\midrule
High interest & -0.071 (0.042) & 0.381*** (0.026) & -0.093*** (0.027)\\
Low interest & 0.009 (0.043) & 0.305*** (0.024) & -0.062* (0.025)\\
\bottomrule
\multicolumn{4}{l}{\textsuperscript{} OLS coefficient on econ\_index with wave FE and standard controls. SE in parentheses.}\\
\multicolumn{4}{l}{*p < .05, **p < .01, ***p < .001.}\\
\end{tabular}

}

\end{table}%

\subsection{E.2 Winner--Loser
Subgroups}\label{e.2-winnerloser-subgroups}

Electoral winners and losers often display different patterns of
democratic support. Table~\ref{tbl-winner} tests whether the
satisfaction--quality decoupling is symmetric across winners (supporters
of the governing party) and losers (supporters of the opposition).

\begin{table}

\caption{\label{tbl-winner}Korea and Taiwan: econ to DV by electoral
winner/loser status}

\centering{

\centering
\begin{threeparttable}
\resizebox{\ifdim\width>\linewidth\linewidth\else\width\fi}{!}{
\begin{tabular}[t]{ll>{\raggedleft\arraybackslash}p{3.5cm}>{\raggedleft\arraybackslash}p{3.5cm}}
\toprule
Country & Group & Dem. always pref. & Sat. w/ democracy\\
\midrule
 & Loser & 0.031 (0.062) & 0.364*** (0.038)\\

\multirow{-2}{*}{\raggedright\arraybackslash Korea} & Winner & -0.052 (0.053) & 0.349*** (0.031)\\
\cmidrule{1-4}
 & Loser & -0.128* (0.059) & 0.406*** (0.033)\\

\multirow{-2}{*}{\raggedright\arraybackslash Taiwan} & Winner & -0.246*** (0.047) & 0.311*** (0.025)\\
\bottomrule
\end{tabular}}
\begin{tablenotes}
\item OLS coefficient on econ\_index with wave FE and standard controls. SE in parentheses. *p < .05, **p < .01, ***p < .001.
\end{tablenotes}
\end{threeparttable}

}

\end{table}%

\section{Appendix F: Authoritarian Rejection --- Full
Results}\label{appendix-f-authoritarian-rejection-full-results}

\subsection{F.1 Pooled Models}\label{f.1-pooled-models}

Table~\ref{tbl-auth-pooled} reports pooled OLS estimates of econ →
authoritarian rejection (index and four component items) for Korea and
Taiwan. The sign test is: if the Korean decoupling reflects democratic
erosion, better economic evaluations should weaken rejection of
authoritarian alternatives. If it reflects only a
performance-satisfaction response, they should strengthen rejection.

\begin{table}

\caption{\label{tbl-auth-pooled}Pooled econ\(\to\)authoritarian
rejection: index and component items}

\centering{

\centering
\begin{tabular}[t]{lllll}
\toprule
\multicolumn{1}{c}{ } & \multicolumn{2}{c}{Korea} & \multicolumn{2}{c}{Taiwan} \\
\cmidrule(l{3pt}r{3pt}){2-3} \cmidrule(l{3pt}r{3pt}){4-5}
Dependent variable & $\beta$ (SE) & N & $\beta$ (SE) & N\\
\midrule
Auth rejection index & -0.082*** (0.018) & 7,519 & -0.040*** (0.011) & 8,729\\
Reject strongman & -0.077** (0.024) & 7,441 & -0.042** (0.015) & 8,452\\
Reject military rule & -0.092*** (0.021) & 7,460 & -0.016 (0.013) & 8,563\\
Reject expert rule & -0.079** (0.026) & 6,307 & -0.019 (0.016) & 6,970\\
Reject single-party & -0.096*** (0.023) & 7,446 & -0.078*** (0.014) & 8,502\\
\bottomrule
\multicolumn{5}{l}{\textsuperscript{} OLS b (econ\_index) with wave FE and standard controls. *p < .05, **p < .01, ***p < .001.}\\
\end{tabular}

}

\end{table}%

\subsection{F.2 Wave-by-Wave Authoritarian
Rejection}\label{f.2-wave-by-wave-authoritarian-rejection}

Table~\ref{tbl-auth-wave} presents the wave-by-wave results for the
authoritarian rejection index and its four component items in both
countries.

\begin{table}

\caption{\label{tbl-auth-wave}Wave-by-wave econ to authoritarian
rejection}

\centering{

\centering\begingroup\fontsize{9}{11}\selectfont

\resizebox{\ifdim\width>\linewidth\linewidth\else\width\fi}{!}{
\begin{tabular}[t]{llrrrrr}
\toprule
Country & Year & Auth rejection index & Reject strongman & Reject military rule & Reject expert rule & Reject single-party\\
\midrule
 & 2003 & -0.046 & -0.077 & 0.091* & -0.113* & -0.080\\

 & 2006 & -0.034 & -0.002 & -0.015 & NA & -0.093*\\

 & 2011 & 0.012 & 0.011 & -0.006 & 0.039 & 0.003\\

 & 2015 & -0.224*** & -0.110 & -0.310*** & -0.227*** & -0.256***\\

 & 2019 & -0.161** & -0.190** & -0.198*** & -0.107 & -0.145*\\

\multirow{-6}{*}{\raggedright\arraybackslash Korea} & 2022 & -0.089 & -0.104 & -0.228*** & 0.023 & -0.048\\
\cmidrule{1-7}
 & 2001 & -0.015 & 0.010 & -0.019 & -0.008 & -0.033\\

 & 2006 & -0.065* & -0.127*** & 0.017 & NA & -0.094**\\

 & 2010 & -0.112*** & -0.113** & -0.070* & -0.096** & -0.173***\\

 & 2014 & -0.088*** & -0.062 & -0.076** & -0.098** & -0.097**\\

 & 2019 & -0.011 & 0.007 & 0.004 & 0.008 & -0.078*\\

\multirow{-6}{*}{\raggedright\arraybackslash Taiwan} & 2022 & 0.074** & 0.063 & 0.069* & 0.113** & 0.042\\
\bottomrule
\multicolumn{7}{l}{\textsuperscript{} OLS coefficient on econ\_index. *p < .05, **p < .01, ***p < .001.}\\
\end{tabular}}
\endgroup{}

}

\end{table}%

\subsection{F.3 Cross-Country
Interaction}\label{f.3-cross-country-interaction}

Table~\ref{tbl-auth-xc} tests whether Korea's positive econ → auth
rejection relationship is significantly different from Taiwan's, using
the same cross-country interaction framework as Appendix D.

\begin{table}

\caption{\label{tbl-auth-xc}Cross-country interaction: econ x Korea on
authoritarian rejection outcomes}

\centering{

\centering\begingroup\fontsize{9}{11}\selectfont

\resizebox{\ifdim\width>\linewidth\linewidth\else\width\fi}{!}{
\begin{tabular}[t]{lrrrrr}
\toprule
Term & Auth rejection index & Reject strongman & Reject military rule & Reject expert rule & Reject single-party\\
\midrule
Econ. eval. & -0.011 (0.012) & -0.016 (0.017) & 0.011 (0.014) & 0.005 (0.018) & -0.036* (0.016)\\
Korea & 0.007 (0.012) & -0.016 (0.016) & 0.042** (0.014) & 0.018 (0.018) & -0.001 (0.015)\\
Econ. x Korea & -0.088*** (0.020) & -0.067* (0.027) & -0.123*** (0.024) & -0.097** (0.030) & -0.071** (0.025)\\
\bottomrule
\multicolumn{6}{l}{\textsuperscript{} OLS. Taiwan = reference. Pooled with wave FE and controls. SE in parentheses. *p < .05, **p < .01, ***p < .001.}\\
\end{tabular}}
\endgroup{}

}

\end{table}%

\section{Appendix G: Item
Specificity}\label{appendix-g-item-specificity}

\subsection{G.1 All Outcomes
Side-by-Side}\label{g.1-all-outcomes-side-by-side}

Table~\ref{tbl-item-spec} presents OLS coefficients for
\texttt{econ\_index} on all nine individual dependent variables for
Korea and Taiwan, plus the Korea--Taiwan gap. The key result is that the
gap between Korea's positive economic effect on the abstract normative
item (``dem always preferable'') and its effect on the authoritarian
rejection items is much larger than the corresponding gap in
Taiwan---confirming that the abstract item is capturing something
distinct from substantive democratic commitment.

\begin{table}

\caption{\label{tbl-item-spec}Item specificity: econ\(\to\)all DVs,
Korea and Taiwan}

\centering{

\centering
\begin{threeparttable}
\begin{tabular}[t]{>{\raggedright\arraybackslash}p{4.5cm}>{\raggedleft\arraybackslash}p{2.8cm}r>{\raggedleft\arraybackslash}p{2.8cm}r}
\toprule
\multicolumn{1}{c}{ } & \multicolumn{2}{c}{Korea} & \multicolumn{2}{c}{Taiwan} \\
\cmidrule(l{3pt}r{3pt}){2-3} \cmidrule(l{3pt}r{3pt}){4-5}
Dependent variable & $\beta$ (SE) & N & $\beta$ (SE) & N\\
\midrule
Satisfaction with democracy & 0.342*** (0.017) & 7,392 & 0.398*** (0.015) & 8,648\\
Satisfaction with government & 0.575*** (0.021) & 7,338 & 0.660*** (0.017) & 8,527\\
Dem always preferable & -0.032 (0.030) & 7,267 & -0.239*** (0.028) & 8,527\\
Democratic extent & 0.059*** (0.006) & 6,655 & 0.114*** (0.008) & 7,795\\
Auth rejection index & -0.082*** (0.018) & 7,519 & -0.040*** (0.011) & 8,729\\
\addlinespace
Reject strongman & -0.077** (0.024) & 7,441 & -0.042** (0.015) & 8,452\\
Reject military rule & -0.092*** (0.021) & 7,460 & -0.016 (0.013) & 8,563\\
Reject expert rule & -0.079** (0.026) & 6,307 & -0.019 (0.016) & 6,970\\
Reject single-party & -0.096*** (0.023) & 7,446 & -0.078*** (0.014) & 8,502\\
\bottomrule
\end{tabular}
\begin{tablenotes}
\item Pooled OLS b (econ\_index) with wave FE and standard controls. SE in parentheses. *p < .05, **p < .01, ***p < .001.
\end{tablenotes}
\end{threeparttable}

}

\end{table}%

\section*{References}\label{references}
\addcontentsline{toc}{section}{References}

\phantomsection\label{refs}
\begin{CSLReferences}{1}{0}
\bibitem[\citeproctext]{ref-Asian-Barometer-Survey2023-uh}
Asian Barometer Survey. (2023). \emph{{Asian Barometer Survey, Waves 1-6
(2001--2021)}} {[}Dataset{]}.

\bibitem[\citeproctext]{ref-Dalton2004-bq}
Dalton, Russell J. (2004). \emph{{Democratic challenges, democratic
choices: The erosion of political support in advanced industrial
democracies}}. Oxford University Press.
\url{https://doi.org/10.1093/acprof:oso/9780199268436.001.0001}

\end{CSLReferences}




\end{document}
